\appendix
\section*{Appendix}

\section{Precedence-Constrained Average-Case Decision Tree} \label{section:PCACDT}

For the average-case objective, bounding the number of tests in a sepcover is insufficient: each test contributes the number of hypotheses that reach it. We therefore compare the total contribution of the selected tests with the corresponding contribution of an optimal sepcover.
We begin with a counting identity.
\begin{observation} \label{obs:dec-tree-averaging}
    If $D$ is a decision tree for $\ProblemPCACDT$, then $\COSTA\br{D}=\sum_{v\in X}\spr{\cH_v}$, where $X$ is the set of all non-leaf vertices of $D$.
\end{observation}
For a subgraph $D'$ of $D$, define $\COSTA\br{D'}=\sum_{v\in V(D')\cap\inner{D}}\spr{\cH_v(D)}$, where $\inner{D}$ is the set of test nodes of $D$. In particular, the last test of a path is counted even though it is a leaf of the path.
Observation~\ref{obs:dec-tree-averaging} allows us to decompose the cost criterion as follows.
\begin{observation} \label{obs:WCdecomposition}
    Let $D$ be any decision tree for an instance of $\ProblemPCACDT$ and let $S=\tests{D}{v}$ for any test node $v$ of $D$, with $D_S$ denoting the path to that node.
    Then,
    \[
    \COSTA\br{D} = \COSTA\br{D_{S}} + \sum_{D'\in D\setminus V\br{D_{S}}} \COSTA\br{D'}.
    \]
%     $$
%     \COSTA\br{D, \cI} = \COSTA\br{S, \cI} + \sum_{D'\in D\setminus S} \COSTA\br{D', \cI}.
%     $$

\end{observation}
\paraTitle{Remark on notation.}
Throughout this subsection, we order $\sepcover{\cI}$ along its path in an optimal average-case decision tree and write $\COSTA\br{\sepcover{\cI}}$ for the cost of that path.
As an immediate corollary, by taking an optimal $D^*$ and the corresponding $\sepcover{\cI}$ we have:
\begin{observation} \label{obs:sepcoverAC}
    Let $\cI$ be any instance of $\ProblemPCACDT$.
    Then, $\COSTA\br{\sepcover{\cI}} \leq \OPT\br{\cI}$.
\end{observation}
We now explore a connection to $\ProblemPCMSSC$, which allows us to use $\COSTA\br{\sepcover{\cI}}$ as a lower bound on $\OPT\br{\cI}$.
% However, since in this case we have an arbitrary probability distrubution on hypotheses, we firstly employ a rounding technique which scales the probabilities to represent them as polynomial-size integers.
% Below we argue that this incurrs only a 2-factor loss in the approximation.
% Formally, for an instance $\cI=\br{\cH, \cT, \preceq, p}$, define $p'(h)=z(h)/s_z$ for each $h\in\cH$, where  $z(h) = \cl{p(h)/q}$ and $s_z=\sum_{h\in\cH}z(h)$, $q=\frac{1}{nm}$, $n=\spr{\cH}$ and $m=\spr{\cT}$.
% We say that $p'$ is a \emph{$q$-rounding} of $p$.
% Note that $\sum_{h\in\cH}p'\br{h}=1$ and thus $p'$ is indeed a probability distribution over $\cH$.
% \begin{lemma} \label{lemma:rounding}
%     Let $\cI=\br{\cH, \cT, \preceq, p}$ be any instance of $\ProblemPCACDT$ and let $p'$ be the $q$-rounding of $p$.
%     Then, for the instance $\cI'=\br{\cH, \cT, \preceq, p'}$ it holds $z(h)\in\bigo(nm)$ and $\OPT\br{\cI}\leq\OPT\br{\cI'} \leq 2 \cdot \OPT\br{\cI}$.
% \end{lemma}
% \begin{proof}
%     Note that the lower bound $\OPT\br{\cI}\leq\OPT\br{\cI'}$ follows from $p\leq p'$, and it remains to upper bound $\OPT\br{\cI'}$.
%     Since $z(h) \leq p(h)/q + 1$, $s_z \leq \frac{1}{q} \sum_{h \in \cH} p(h) + n = nm + n$.
%     Let $D'$ be an optimal decision tree for $\cI'$ and let $D$ be the optimal decision tree for $\cI$.
%     By $p\leq p'$ and the optimality of $D'$ for $\cI'$, we get $\COSTA\br{D', \cI} \leq \COSTA\br{D', \cI'}\leq \COSTA\br{D, \cI'}$.
%     By Observation~\ref{obs:dec-tree-averaging}, $\COSTA\br{D, \cI'}=\sum_{v\in V(D)}p'(v)\cdot\spr{\cH_v}$.
%     Hence, by $s_z p'(h)\leq1+p(h)/q$,
%     \begin{align*}
%     \COSTA\br{D', \cI} & \leq  \frac{1}{qs_z}\cdot\sum_{v\in V(D)}p(v)\cdot\spr{\cH_v} + \frac{1}{s_z}\sum_{v\in V(D)} \spr{\cH_v}
%     \\&\leq
%      \frac{n m}{s_z}\cdot \COSTA\br{\br{D, \cI} + 1}
%     \\&\leq
%     \COSTA\br{D, \cI} + 1
%     \leq
%     2 \cdot \COSTA\br{D, \cI}
%     \end{align*}
%     where the second inequality is by Observation~\ref{obs:dec-tree-averaging} and $\spr{\cH_v}\leq n$ for each node $v$, the third inequality is by $s_z\geq nm$ and the last inequality follows since $\COSTA\br{D, \cI} \geq \sum_{h\in\cH}p\br{h} = 1$.
% %     Let $q=\frac{1}{nm}$. For every $h\in\cH$, let $p'(h) = \cl{p(h)/q}$. Therefore, $p'(h) \leq p(h)/q + 1$ and as a consequence, $p'\br{\cH} \leq \sum_{h \in \cH} \br {\frac{p(h)}{q} + 1} = \frac{1}{q} \sum_{h \in \cH} p(h) + n = nm + n = \bigo(nm)$.
% %
% %     We now argue that $\OPT\br{\cI'} \leq 2 \cdot \OPT\br{\cI}$. Let $D'$ be an optimal decision tree for $\cI'$ and let $D$ be the optimal decision tree for $\cI$.
% %     \begin{align*}
% %     \COSTA\br{D', \cI} &\leq q\cdot \COSTA\br{D', \cI'}\\&\leq
% %     q\cdot \COSTA\br{D, \cI'}
% %     \\&\leq
% %     \COSTA\br{D, \cI} + q\cdot\sum_{h\in\cH} \COST\br{D, \cI', h}
% %     \\&\leq
% %     \COSTA\br{D, \cI} + q\cdot n\cdot m
% %     \\&=
% %     \COSTA\br{D, \cI} + 1
% %     \\&\leq
% %     2 \cdot \COSTA\br{D, \cI}
% %     \end{align*}
% %     where the first inequality follows since clearly, for every $h\in\cH$, we have that $p(u)\leq q\cdot p'(h)$, the second inequality is due to the optimality of $D$ for $\cI$, the third inequality follows since for every $h\in\cH$, $q\cdot p'(h) \leq p(h) + q$, the fourth inequality follows since the depth of $D$ is at most $m$, the equality is by definition of $q$ and the last inequality follows since $\COSTA\br{D, \cI} \geq p\br{\cH} = 1$.
% \end{proof}
% Therefore, from now on we will assume that the input instance $\cI=\br{\cH, \cT, \preceq, p}$ of $\ProblemPCACDT$ is such that for each $h\in\cH$, $p(h)=i/j$ for integers $i<j=\bigo(nm)$, without impacting the asymptotic approximation ratio of our algorithm.
Hence, we have an analogue of Lemma \ref{lemma:cover_sep_worst_case}.
\begin{lemma}\label{lemma:cover_sep_average_case}
    Let $\cI=\br{\cH,\cT,\preceq}$ be a $\ProblemPCACDT$ instance, and let $\cC^*$ be an optimal solution to $\ProblemPCMSSC$ on $\br{\cU_{\xi},\xi\br{\cT},\preceq_{\xi},f=1/4}$.
    %, where a test $t$ covers element $u \in \cH$ if for $u\in U_{t,j}$, $\spr{U_{t,j}}\leq \frac{3}{4}\cdot \spr{\cH}$.
    Then, $\coverageTime\br{\cC^*} \leq \COSTA\br{\sepcover{\cI}}$.
\end{lemma}
\begin{proof}
    Let $\sigma$ be the shortest prefix of $\xi(\sepcover{\cI})$ that covers at least $\spr{\cU_{\xi}}/4$ hypotheses. By the proof of Lemma~\ref{lemma:cover_sep_worst_case}, this prefix exists and is precedence-closed. For $h\in\cU_{\xi}$, let $d(h)$ be the number of tests on the sepcover path performed when the target is $h$. If $h$ leaves the path by the end of $\sigma$, the same proof shows that $h$ is covered no later than its departure, so $\coverageTime\br{h,\sigma}\leq d(h)$. Otherwise, $\coverageTime\br{h,\sigma}\leq\spr{\sigma}\leq d(h)$. Summing gives
    \[
    \coverageTime\br{\cC^*}\leq\coverageTime\br{\sigma}
    =\sum_{h\in\cU_{\xi}}\coverageTime\br{h,\sigma}
    \leq\sum_{h\in\cU_{\xi}}d(h)
    \leq\COSTA\br{\sepcover{\cI}}.\qedhere
    \]
\end{proof}

The following is analogous to Observation~\ref{obs:DS} but requires a more careful consideration.
\begin{observation} \label{obs:DS-average}
    Consider $D_{\cC}$ induced by Algorithm~\ref{alg:worstDecisionTree}, which uses an $\br{\alpha,\beta}$-approximation algorithm for $\ProblemPCMSSC$ as the $\blackBox$ subroutine.
    Then, $\COSTA\br{D_{\cC}} \leq 2\cdot\alpha \cdot \COSTA\br{\sepcover{\cI}}$, where $\cI=\br{\cH,\cT,\preceq}$ is the input instance to $\ProblemPCACDT$.
\end{observation}
\begin{proof}
    Let $\cC_0$ be the shortest prefix of $\cC$ covering at least $\spr{\cU_{\xi}}/(4\cdot\beta)$ hypotheses. By the argument preceding Algorithm~\ref{alg:worstDecisionTree}, the path $D_{\cC}$ ends by the last test associated with $\cC_0$. If $h\in\cU_{\xi}$ is first covered by $\cC_0$ at position $j$, then $h$ belongs to a reply of $t_j$ of size at most $3\cdot\spr{\cH}/4\leq (5\cdot\beta-1)\cdot\spr{\cH}/(5\cdot\beta)$. Thus $h$ leaves the current path by test $t_j$, unless the path has already ended. Its contribution to $\COSTA\br{D_{\cC}}$ is at most $j=\coverageTime\br{h,\cC_0}$. If $h$ is not covered by $\cC_0$, its contribution is at most $\spr{\cC_0}=\coverageTime\br{h,\cC_0}$.
    There is at most one hypothesis in $\cH\setminus\cU_{\xi}$, and its contribution to $\COSTA\br{D_{\cC}}$ is at most $\spr{\cC_0}$. The last set of $\cC_0$ covers a new hypothesis, so $\spr{\cC_0}\leq\coverageTime\br{\cC_0}$. Moreover, $\coverageTime\br{h,\cC_0}=\min\brc{\coverageTime\br{h,\cC},\spr{\cC_0}}$ for every $h\in\cU_{\xi}$, so $\coverageTime\br{\cC_0}\leq\coverageTime\br{\cC}$. The approximation guarantee for $\blackBox$ and Lemma~\ref{lemma:cover_sep_average_case} give
    \[
    \COSTA\br{D_{\cC}}
    \leq \coverageTime\br{\cC_0}+\spr{\cC_0}
    \leq 2\cdot\coverageTime\br{\cC_0}
    \leq 2\cdot\coverageTime\br{\cC}
    \leq 2\cdot\alpha\cdot\COSTA\br{\sepcover{\cI}}. \qedhere
    \]
\end{proof}

\begin{theorem}\label{thm:PCMSSC-to-PCACDT}
    \setlength{\emergencystretch}{2em}
    If there is a polynomial-time bicriteria $\br{\alpha, \beta}$-approximation algorithm for $\ProblemPCMSSC$, then there is a polynomial-time $\bigo\br{\frac{\alpha}{\log\br{\frac{5\cdot\beta}{5\cdot\beta -1}}}\cdot \log n}$-approximation algorithm for $\ProblemPCACDT$.
    % In particular when $\mu = \bigo\br{1}$, the ratio is $\bigo\br{\beta \cdot \log \br{m+n}}$.
\end{theorem}
\begin{proof}
    \setlength{\emergencystretch}{2em}
    Let $D$ be the decision tree induced by \ProcDecisionTree on $\cI=\br{\cH,\cT,\preceq}$, using a $\br{\alpha,\beta}$-approximation for $\ProblemPCMSSC$ as $\blackBox$.
    We argue by induction on $n$ that
    \begin{equation} \label{eq:D-average}
     \COSTA\br{D} \leq \frac{2\cdot\alpha}{\log\br{\frac{5\cdot\beta}{5\cdot\beta -1}}}\cdot \log n \cdot \OPT\br{\cI}.
    \end{equation}
    For $n=1$, $\COSTA\br{D}=0$. Assume that \eqref{eq:D-average} holds for every recursive call on $\cI'=\br{\cH',\cT,\preceq}$. Then
\begin{align*}
    &\COSTA\br{D}=\COSTA\br{D_{\cC}}+\sum_{D'\in D\setminus V\br{D_{\cC}}}\COSTA\br{D'} &&\\
    &\leq 2\cdot\alpha\cdot\COSTA\br{\sepcover{\cI}}+\sum_{D'\in D\setminus V\br{D_{\cC}}}\frac{2\cdot\alpha}{\log\br{\frac{5\cdot\beta}{5\cdot\beta-1}}}\cdot\log\spr{\cH'}\cdot\OPT\br{\cI'} &&\\
    &\leq 2\cdot\alpha\cdot\OPT\br{\cI}+\frac{2\cdot\alpha}{\log\br{\frac{5\cdot\beta}{5\cdot\beta-1}}}\cdot\log\br{\frac{5\cdot\beta-1}{5\cdot\beta}\cdot n}\cdot\OPT\br{\cI} &&\\
    &=2\cdot\alpha\cdot\OPT\br{\cI}+\frac{2\cdot\alpha}{\log\br{\frac{5\cdot\beta}{5\cdot\beta-1}}}\cdot\log n\cdot\OPT\br{\cI}-2\cdot\alpha\cdot\OPT\br{\cI} &&\\
    &=\frac{2\cdot\alpha}{\log\br{\frac{5\cdot\beta}{5\cdot\beta-1}}}\cdot\log n\cdot\OPT\br{\cI}. &&
\end{align*}
The first equality uses Observation~\ref{obs:WCdecomposition}. The first inequality uses the induction hypothesis and Observation~\ref{obs:DS-average}. For the second one, use Observation~\ref{obs:sepcoverAC}, $\spr{\cH'}\leq\br{1-\frac45\cdot\frac{1}{4\cdot\beta}}\cdot n=\frac{5\cdot\beta-1}{5\cdot\beta}\cdot n$, and Lemma~\ref{lemma:subinstances_avg}, since the hypothesis sets of the recursive calls are disjoint. The next equality follows by expanding the logarithm as in the worst-case proof.
\end{proof}
